\documentclass[12pt]{article}
\usepackage[utf8]{inputenc}
\usepackage[spanish]{babel}
\usepackage{amsmath, amssymb, amsthm}
\usepackage{enumitem}
\usepackage[margin=2cm]{geometry}
\usepackage{float}
\usepackage{parskip}


\begin{document}

\section*{Problemset 2 -- Soluciones}

%==============================================================
\section{Problema 1: Variable aleatoria discreta}

Se presenta una tabla de distribución de probabilidad de una variable aleatoria discreta $X$. Se pide:

\begin{enumerate}[label=\alph*)]
    \item Describir la función de masa de probabilidad (FMP) de $X$ y representarla gráficamente.
    \item Describir la función de distribución acumulada (FDA) de $X$ y representarla gráficamente.
    \item Calcular la probabilidad indicada.
    \item Calcular $E(X)$ y $\operatorname{Var}(X)$.
    \item Comparar los resultados obtenidos para distintas variables.
    \item Describir la distribución de probabilidades de $X$ y comparar.
\end{enumerate}

\subsection{Solución}

\textbf{a)} La función de masa de probabilidad (FMP) de una variable aleatoria discreta $X$ se define como
$$p(x) = P(X = x), \quad \text{para cada valor posible } x,$$
con las condiciones $p(x)\geq 0$ para todo $x$ y $\displaystyle\sum_x p(x) = 1$.

Se representa mediante una gráfica lineal: sobre cada valor $x$ con $p(x)>0$ se coloca un punto a altura $p(x)$.

\textbf{b)} La función de distribución acumulada se define como
$$F(x) = P(X \leq x) = \sum_{t \leq x} p(t).$$
Es una función escalonada, no decreciente, con $\lim_{x\to -\infty}F(x)=0$ y $\lim_{x\to\infty}F(x)=1$.

\textbf{d)} El valor esperado y la varianza se calculan como
$$E(X) = \sum_x x\, p(x),$$
$$\operatorname{Var}(X) = E(X^2) - [E(X)]^2 = \sum_x x^2\, p(x) - \left[\sum_x x\, p(x)\right]^2.$$

\textbf{e) y f)} Se comparan los resultados de $E(X)$, $\operatorname{Var}(X)$ y la forma de la distribución entre las distintas variables o escenarios planteados, verificando que las sumas de probabilidades sean 1 y que las propiedades se cumplan.

%==============================================================
\section{Problema 2: Variable aleatoria continua}

Sea $X$ una variable aleatoria continua con función de densidad de probabilidad dada. Se pide:

\begin{enumerate}[label=\alph*)]
    \item Describir la función de densidad de probabilidad (FDP) de $X$ y representarla gráficamente.
    \item Describir la función de distribución acumulada (FDA) de $X$ y representarla gráficamente.
    \item Calcular la probabilidad indicada.
    \item Calcular $E(X)$ y $\operatorname{Var}(X)$.
    \item Comparar con los resultados de otra variable $N$.
    \item Describir la distribución de $X$ y comparar con $N$.
\end{enumerate}

\subsection{Solución}

\textbf{a)} La función de densidad de probabilidad $f(x)$ debe satisfacer:
$$f(x) \geq 0 \quad \forall\, x, \qquad \int_{-\infty}^{\infty} f(x)\,dx = 1.$$
Para cualquier intervalo $[a,b]$:
$$P(a \leq X \leq b) = \int_a^b f(x)\,dx.$$

\textbf{b)} La FDA se obtiene integrando la FDP:
$$F(x) = P(X \leq x) = \int_{-\infty}^{x} f(t)\,dt.$$
Es una función continua, no decreciente, con $F(-\infty)=0$ y $F(+\infty)=1$.

\textbf{d)} Para una variable continua:
$$E(X) = \int_{-\infty}^{\infty} x\, f(x)\,dx,$$
$$E(X^2) = \int_{-\infty}^{\infty} x^2\, f(x)\,dx,$$
$$\operatorname{Var}(X) = E(X^2) - [E(X)]^2.$$

\textbf{e) y f)} Se comparan $E(X)$, $\operatorname{Var}(X)$ y la forma de la distribución de $X$ con los de $N$, verificando si tienen la misma estructura distribucional.

%==============================================================
\section{Problema 3: Distribución y comparación}

Se plantea un problema sobre una distribución de probabilidad específica y se pide comparar con otra distribución o variable.

\subsection{Solución}

Se identifican los parámetros de la distribución, se calculan $E(X)$, $\operatorname{Var}(X)$ y se comparan con los de la otra variable o distribución indicada, verificando si las propiedades distribucionales coinciden.

%==============================================================
\section{Problema 4: Intervalos de confianza}

Se plantea un problema sobre la construcción de un intervalo de confianza para un parámetro poblacional.

\subsection{Solución}

Dependiendo del contexto (media, proporción, varianza), se construye el intervalo de confianza apropiado:

Para una media $\mu$ con $\sigma$ conocida:
$$\bar{x} \pm z_{\alpha/2}\,\frac{\sigma}{\sqrt{n}}.$$

Para una media $\mu$ con $\sigma$ desconocida:
$$\bar{x} \pm t_{\alpha/2,\, n-1}\,\frac{s}{\sqrt{n}}.$$

Para una proporción $p$:
$$\hat{p} \pm z_{\alpha/2}\,\sqrt{\frac{\hat{p}(1-\hat{p})}{n}}.$$

%==============================================================
\section{Problema 5: Comparación de dos poblaciones}

Se plantea un problema de comparación de dos poblaciones o de prueba de hipótesis sobre la diferencia de medias o proporciones.

\subsection{Solución}

Para comparar dos medias $\mu_1$ y $\mu_2$ con muestras independientes:

\textbf{Estadístico de prueba (muestras grandes):}
$$Z = \frac{\bar{X} - \bar{Y}}{\sqrt{\dfrac{S_1^2}{m} + \dfrac{S_2^2}{n}}}.$$

\textbf{Prueba $t$ con dos muestras (varianzas desconocidas):}
$$T = \frac{\bar{X} - \bar{Y}}{\sqrt{\dfrac{S_1^2}{m} + \dfrac{S_2^2}{n}}},$$
con grados de libertad aproximados dados por la fórmula de Welch--Satterthwaite.

\textbf{Intervalo de confianza para $\mu_1 - \mu_2$:}
$$(\bar{x} - \bar{y}) \pm t_{\alpha/2,\,\nu}\,\sqrt{\frac{s_1^2}{m} + \frac{s_2^2}{n}},$$
donde $\nu$ son los grados de libertad aproximados.

\textbf{Para datos apareados:} se trabaja con las diferencias $d_i = x_i - y_i$ y se aplica una prueba $t$ de una muestra a las diferencias:
$$T = \frac{\bar{d}}{s_d / \sqrt{n}},$$
con $n-1$ grados de libertad.

\section{Problema 6: Estrellas variables}

En un cúmulo estelar se han identificado $N=14$ estrellas variables, de las cuales $10$ son de tipo Cefeidas y $4$ son de tipo RR~Lyrae. Se seleccionan $n=5$ estrellas al azar (sin reemplazo) para un estudio.

\begin{enumerate}[label=\alph*)]
    \item ¿Cuál es la media (valor esperado) del número de estrellas de tipo RR~Lyrae en la muestra?
    \item Determine la probabilidad de que el número de Cefeidas entre las $5$ estrellas observadas exceda su valor medio en más de una desviación estándar.
    \item Suponga que ahora se tiene un conjunto de $300$ estrellas variables, de las cuales $30$ son de tipo Cefeidas. Si se escogen $4$ estrellas al azar, ¿cuál es la probabilidad de que al menos $1$ estrella sea Cefeidas? Analice la posibilidad de utilizar una distribución diferente a la hipergeométrica y explique.
\end{enumerate}

\subsection{Solución}

El experimento consiste en extraer $n$ estrellas sin reemplazo de una población finita de $N$ estrellas, donde $M$ son de un tipo de interés. Esto corresponde a una distribución hipergeométrica. Según Devore (sección 3.5), si $X$ es el número de ``éxitos'' en la muestra, entonces:
$$X \sim \text{Hipergeométrica}(n, M, N), \qquad P(X=x)=h(x;n,M,N)=\frac{\binom{M}{x}\binom{N-M}{n-x}}{\binom{N}{n}}$$

\subsection*{Inciso a)}

Sea $X_{RR}$ el número de estrellas RR~Lyrae en la muestra de $n=5$. Aquí $N=14$, $M=4$ (RR~Lyrae), $n=5$.

La esperanza de una variable hipergeométrica es (Devore, sección 3.5):
$$E(X_{RR})=n\cdot\frac{M}{N}=5\cdot\frac{4}{14}=\frac{20}{14}=\frac{10}{7}\approx 1.4286$$

Por tanto, en promedio se esperan aproximadamente $1.43$ estrellas RR~Lyrae en la muestra de $5$.

\subsection*{Inciso b)}

Ahora sea $X_{C}$ el número de Cefeidas en la muestra. Aquí $N=14$, $M=10$ (Cefeidas), $n=5$.

\textbf{Valor esperado:}
$$E(X_C)=n\cdot\frac{M}{N}=5\cdot\frac{10}{14}=\frac{50}{14}=\frac{25}{7}\approx 3.5714$$

\textbf{Varianza:} La varianza de la hipergeométrica incluye el factor de corrección por población finita:
$$V(X_C)=n\cdot\frac{M}{N}\cdot\left(1-\frac{M}{N}\right)\cdot\frac{N-n}{N-1}$$
$$V(X_C)=5\cdot\frac{10}{14}\cdot\frac{4}{14}\cdot\frac{9}{13}=\frac{5\cdot 10\cdot 4\cdot 9}{14\cdot 14\cdot 13}=\frac{1800}{2548}\approx 0.7064$$

La desviación estándar es:
$$\sigma_{X_C}=\sqrt{0.7064}\approx 0.8405$$

Se pide:
$$P\!\left(X_C > E(X_C)+\sigma_{X_C}\right)=P(X_C > 3.5714+0.8405)=P(X_C>4.412)$$

Como $X_C$ solo toma valores enteros $\{0,1,2,3,4,5\}$, se tiene:
$$P(X_C>4.412)=P(X_C=5)$$

Calculamos:
$$P(X_C=5)=\frac{\binom{10}{5}\binom{4}{0}}{\binom{14}{5}}=\frac{252\cdot 1}{2002}=\frac{252}{2002}\approx 0.1259$$

Por tanto, la probabilidad de que el número de Cefeidas exceda su valor medio en más de una desviación estándar es aproximadamente $0.126$.

\subsection*{Inciso c)}

Ahora $N=300$, $M=30$ (Cefeidas), $n=4$. Sea $Y$ el número de Cefeidas en la muestra.

\textbf{Cálculo exacto (hipergeométrica):}
$$P(Y\geq 1)=1-P(Y=0)=1-\frac{\binom{270}{4}}{\binom{300}{4}}$$

$$P(Y=0)=\frac{\binom{270}{4}}{\binom{300}{4}}=\frac{270\cdot 269\cdot 268\cdot 267}{300\cdot 299\cdot 298\cdot 297}$$

$$P(Y=0)=\frac{270}{300}\cdot\frac{269}{299}\cdot\frac{268}{298}\cdot\frac{267}{297}\approx 0.9\cdot 0.8997\cdot 0.8993\cdot 0.8990\approx 0.6520$$

$$P(Y\geq 1)\approx 1-0.6520=0.3480$$

\textbf{Aproximación binomial:} Dado que $n/N=4/300\approx 0.013<0.05$, se cumple la regla empírica (Devore, sección 3.5) que permite aproximar la hipergeométrica por una binomial con $p=M/N=30/300=0.1$:
$$Y\approx \text{Bin}(4,\,0.1)$$
$$P(Y\geq 1)=1-P(Y=0)=1-(0.9)^4=1-0.6561=0.3439$$

\textbf{Análisis:} La aproximación binomial ($0.3439$) es muy cercana al valor exacto hipergeométrico ($0.3480$), con una diferencia menor a $0.005$. Esto es consistente con la regla empírica: cuando $n/N<0.05$, el muestreo sin reemplazo se comporta casi como muestreo con reemplazo, y la distribución binomial es una aproximación adecuada. La ventaja de la aproximación binomial es su simplicidad computacional, ya que no requiere calcular coeficientes binomiales grandes.

\end{document}